% Options for packages loaded elsewhere
% Options for packages loaded elsewhere
\PassOptionsToPackage{unicode}{hyperref}
\PassOptionsToPackage{hyphens}{url}
%
\documentclass[
  ignorenonframetext,
  aspectratio=169,
  spanish,
]{beamer}
\newif\ifbibliography
\usepackage{pgfpages}
\setbeamertemplate{caption}[numbered]
\setbeamertemplate{caption label separator}{: }
\setbeamercolor{caption name}{fg=normal text.fg}
\beamertemplatenavigationsymbolshorizontal
% remove section numbering
\setbeamertemplate{part page}{
  \centering
  \begin{beamercolorbox}[sep=16pt,center]{part title}
    \usebeamerfont{part title}\insertpart\par
  \end{beamercolorbox}
}
\setbeamertemplate{section page}{
  \centering
  \begin{beamercolorbox}[sep=12pt,center]{section title}
    \usebeamerfont{section title}\insertsection\par
  \end{beamercolorbox}
}
\setbeamertemplate{subsection page}{
  \centering
  \begin{beamercolorbox}[sep=8pt,center]{subsection title}
    \usebeamerfont{subsection title}\insertsubsection\par
  \end{beamercolorbox}
}
% Prevent slide breaks in the middle of a paragraph
\widowpenalties 1 10000
\raggedbottom
\AtBeginPart{
  \frame{\partpage}
}
\AtBeginSection{
  \ifbibliography
  \else
    \frame{\sectionpage}
  \fi
}
\AtBeginSubsection{
  \frame{\subsectionpage}
}
\usepackage{iftex}
\ifPDFTeX
  \usepackage[T1]{fontenc}
  \usepackage[utf8]{inputenc}
  \usepackage{textcomp} % provide euro and other symbols
\else % if luatex or xetex
  \usepackage{unicode-math} % this also loads fontspec
  \defaultfontfeatures{Scale=MatchLowercase}
  \defaultfontfeatures[\rmfamily]{Ligatures=TeX,Scale=1}
\fi
\usepackage{lmodern}

\usetheme[]{default}
\usefonttheme[]{default}
\ifPDFTeX\else
  % xetex/luatex font selection
\fi
% Use upquote if available, for straight quotes in verbatim environments
\IfFileExists{upquote.sty}{\usepackage{upquote}}{}
\IfFileExists{microtype.sty}{% use microtype if available
  \usepackage[]{microtype}
  \UseMicrotypeSet[protrusion]{basicmath} % disable protrusion for tt fonts
}{}
\makeatletter
\@ifundefined{KOMAClassName}{% if non-KOMA class
  \IfFileExists{parskip.sty}{%
    \usepackage{parskip}
  }{% else
    \setlength{\parindent}{0pt}
    \setlength{\parskip}{6pt plus 2pt minus 1pt}}
}{% if KOMA class
  \KOMAoptions{parskip=half}}
\makeatother

\usepackage{color}
\usepackage{fancyvrb}
\newcommand{\VerbBar}{|}
\newcommand{\VERB}{\Verb[commandchars=\\\{\}]}
\DefineVerbatimEnvironment{Highlighting}{Verbatim}{commandchars=\\\{\}}
% Add ',fontsize=\small' for more characters per line
\usepackage{framed}
\definecolor{shadecolor}{RGB}{241,243,245}
\newenvironment{Shaded}{\begin{snugshade}}{\end{snugshade}}
\newcommand{\AlertTok}[1]{\textcolor[rgb]{0.68,0.00,0.00}{#1}}
\newcommand{\AnnotationTok}[1]{\textcolor[rgb]{0.37,0.37,0.37}{#1}}
\newcommand{\AttributeTok}[1]{\textcolor[rgb]{0.40,0.45,0.13}{#1}}
\newcommand{\BaseNTok}[1]{\textcolor[rgb]{0.68,0.00,0.00}{#1}}
\newcommand{\BuiltInTok}[1]{\textcolor[rgb]{0.00,0.23,0.31}{#1}}
\newcommand{\CharTok}[1]{\textcolor[rgb]{0.13,0.47,0.30}{#1}}
\newcommand{\CommentTok}[1]{\textcolor[rgb]{0.37,0.37,0.37}{#1}}
\newcommand{\CommentVarTok}[1]{\textcolor[rgb]{0.37,0.37,0.37}{\textit{#1}}}
\newcommand{\ConstantTok}[1]{\textcolor[rgb]{0.56,0.35,0.01}{#1}}
\newcommand{\ControlFlowTok}[1]{\textcolor[rgb]{0.00,0.23,0.31}{\textbf{#1}}}
\newcommand{\DataTypeTok}[1]{\textcolor[rgb]{0.68,0.00,0.00}{#1}}
\newcommand{\DecValTok}[1]{\textcolor[rgb]{0.68,0.00,0.00}{#1}}
\newcommand{\DocumentationTok}[1]{\textcolor[rgb]{0.37,0.37,0.37}{\textit{#1}}}
\newcommand{\ErrorTok}[1]{\textcolor[rgb]{0.68,0.00,0.00}{#1}}
\newcommand{\ExtensionTok}[1]{\textcolor[rgb]{0.00,0.23,0.31}{#1}}
\newcommand{\FloatTok}[1]{\textcolor[rgb]{0.68,0.00,0.00}{#1}}
\newcommand{\FunctionTok}[1]{\textcolor[rgb]{0.28,0.35,0.67}{#1}}
\newcommand{\ImportTok}[1]{\textcolor[rgb]{0.00,0.46,0.62}{#1}}
\newcommand{\InformationTok}[1]{\textcolor[rgb]{0.37,0.37,0.37}{#1}}
\newcommand{\KeywordTok}[1]{\textcolor[rgb]{0.00,0.23,0.31}{\textbf{#1}}}
\newcommand{\NormalTok}[1]{\textcolor[rgb]{0.00,0.23,0.31}{#1}}
\newcommand{\OperatorTok}[1]{\textcolor[rgb]{0.37,0.37,0.37}{#1}}
\newcommand{\OtherTok}[1]{\textcolor[rgb]{0.00,0.23,0.31}{#1}}
\newcommand{\PreprocessorTok}[1]{\textcolor[rgb]{0.68,0.00,0.00}{#1}}
\newcommand{\RegionMarkerTok}[1]{\textcolor[rgb]{0.00,0.23,0.31}{#1}}
\newcommand{\SpecialCharTok}[1]{\textcolor[rgb]{0.37,0.37,0.37}{#1}}
\newcommand{\SpecialStringTok}[1]{\textcolor[rgb]{0.13,0.47,0.30}{#1}}
\newcommand{\StringTok}[1]{\textcolor[rgb]{0.13,0.47,0.30}{#1}}
\newcommand{\VariableTok}[1]{\textcolor[rgb]{0.07,0.07,0.07}{#1}}
\newcommand{\VerbatimStringTok}[1]{\textcolor[rgb]{0.13,0.47,0.30}{#1}}
\newcommand{\WarningTok}[1]{\textcolor[rgb]{0.37,0.37,0.37}{\textit{#1}}}

\usepackage{longtable,booktabs,array}
\usepackage{calc} % for calculating minipage widths
\usepackage{caption}
% Make caption package work with longtable
\makeatletter
\def\fnum@table{\tablename~\thetable}
\makeatother
\usepackage{graphicx}
\makeatletter
\newsavebox\pandoc@box
\newcommand*\pandocbounded[1]{% scales image to fit in text height/width
  \sbox\pandoc@box{#1}%
  \Gscale@div\@tempa{\textheight}{\dimexpr\ht\pandoc@box+\dp\pandoc@box\relax}%
  \Gscale@div\@tempb{\linewidth}{\wd\pandoc@box}%
  \ifdim\@tempb\p@<\@tempa\p@\let\@tempa\@tempb\fi% select the smaller of both
  \ifdim\@tempa\p@<\p@\scalebox{\@tempa}{\usebox\pandoc@box}%
  \else\usebox{\pandoc@box}%
  \fi%
}
% Set default figure placement to htbp
\def\fps@figure{htbp}
\makeatother



\ifLuaTeX
\usepackage[bidi=basic]{babel}
\else
\usepackage[bidi=default]{babel}
\fi
% get rid of language-specific shorthands (see #6817):
\let\LanguageShortHands\languageshorthands
\def\languageshorthands#1{}


\setlength{\emergencystretch}{3em} % prevent overfull lines

\providecommand{\tightlist}{%
  \setlength{\itemsep}{0pt}\setlength{\parskip}{0pt}}



 


\setbeamertemplate{headline}{} % Elimina la cabecera del tema Madrid
\setbeamertemplate{section in head/foot}{} % Limpia la numeración de secciones
\renewcommand{\contentsname}{Índice}
\renewcommand{\refname}{Referencias}
\renewcommand{\figurename}{Figura}
\renewcommand{\tablename}{Tabla}
\logo{%
  \begin{tabular}{c}
    \includegraphics[height=1.2cm]{EscudoU.png}
  \end{tabular}%
}
% --- CONFIGURACIÓN DE COLORES INSTITUCIONALES ---
\definecolor{UnicorGreen}{HTML}{024426}
\definecolor{UnicorText}{HTML}{F9F9F9}
\definecolor{UnicorWhite}{HTML}{FFFFFF}
\definecolor{UnicorGold}{HTML}{D1B061}

% Aplicación al fondo y texto general
\setbeamercolor{background canvas}{bg=UnicorGreen}
\setbeamercolor{normal text}{fg=UnicorText}

% Aplicación a títulos y cabeceras
%\setbeamercolor{title}{fg=UnicorWhite}
\setbeamercolor{frametitle}{fg=UnicorWhite}
\setbeamercolor{framesubtitle}{fg=UnicorWhite}

% Aplicación a viñetas y estructura
\setbeamercolor{structure}{fg=UnicorWhite}
\setbeamercolor{itemize item}{fg=UnicorWhite}
\setbeamercolor{itemize subitem}{fg=UnicorWhite}
\setbeamercolor{enumerate item}{fg=UnicorWhite}

% Forzar a LaTeX a repintar el texto general al iniciar el documento
\AtBeginDocument{\usebeamercolor[fg]{normal text}}

% Aplicación a hipervínculos
\hypersetup{colorlinks=true, urlcolor=UnicorGold, linkcolor=UnicorWhite}

% --- ACTIVAR ESTILO DE BLOQUES Y SOMBRAS (ESTILO MADRID) ---
\useinnertheme[shadow=true]{rounded}

% --- COLORES DEL RECUADRO DEL TÍTULO PRINCIPAL ---
% Fondo dorado con texto verde
\setbeamercolor{title}{bg=UnicorGold, fg=UnicorGreen}
\setbeamercolor{frametitle}{fg=UnicorWhite} % Mantiene el título de las diapositivas normal
\setbeamercolor{framesubtitle}{fg=UnicorWhite}

% --- COLORES PARA TEOREMAS, DEFINICIONES Y BLOQUES ---
% Título del bloque (Fondo dorado, texto verde)
\setbeamercolor{block title}{bg=UnicorGold, fg=UnicorGreen}
% Cuerpo del bloque (Fondo blanco, texto verde)
\setbeamercolor{block body}{bg=UnicorWhite, fg=UnicorGreen}

\AtBeginDocument{
  \author{%
    \begin{tabular}{c}
      Mario A. Morales Rivera \\
      \footnotesize \texttt{mamorales@correo.unicordoba.edu.co} \\
      \footnotesize Universidad de Córdoba
    \end{tabular}
  }
  \institute{}
}
% --- LIMPIAR DIAPOSITIVAS DE SECCIÓN (ELIMINAR "SECTION X") ---
\setbeamertemplate{section page}{
  \begin{center}
    \begin{beamercolorbox}[sep=12pt,center,rounded=true,shadow=true]{title}
      \usebeamerfont{title}\insertsection\par
    \end{beamercolorbox}
  \end{center}
}
\makeatletter
\@ifpackageloaded{tcolorbox}{}{\usepackage[skins,breakable]{tcolorbox}}
\@ifpackageloaded{fontawesome5}{}{\usepackage{fontawesome5}}
\definecolor{quarto-callout-color}{HTML}{909090}
\definecolor{quarto-callout-note-color}{HTML}{0758E5}
\definecolor{quarto-callout-important-color}{HTML}{CC1914}
\definecolor{quarto-callout-warning-color}{HTML}{EB9113}
\definecolor{quarto-callout-tip-color}{HTML}{00A047}
\definecolor{quarto-callout-caution-color}{HTML}{FC5300}
\definecolor{quarto-callout-color-frame}{HTML}{acacac}
\definecolor{quarto-callout-note-color-frame}{HTML}{4582ec}
\definecolor{quarto-callout-important-color-frame}{HTML}{d9534f}
\definecolor{quarto-callout-warning-color-frame}{HTML}{f0ad4e}
\definecolor{quarto-callout-tip-color-frame}{HTML}{02b875}
\definecolor{quarto-callout-caution-color-frame}{HTML}{fd7e14}
\makeatother
\makeatletter
\@ifpackageloaded{caption}{}{\usepackage{caption}}
\AtBeginDocument{%
\ifdefined\contentsname
  \renewcommand*\contentsname{Tabla de contenidos}
\else
  \newcommand\contentsname{Tabla de contenidos}
\fi
\ifdefined\listfigurename
  \renewcommand*\listfigurename{Listado de Figuras}
\else
  \newcommand\listfigurename{Listado de Figuras}
\fi
\ifdefined\listtablename
  \renewcommand*\listtablename{Listado de Tablas}
\else
  \newcommand\listtablename{Listado de Tablas}
\fi
\ifdefined\figurename
  \renewcommand*\figurename{Figura}
\else
  \newcommand\figurename{Figura}
\fi
\ifdefined\tablename
  \renewcommand*\tablename{Tabla}
\else
  \newcommand\tablename{Tabla}
\fi
}
\@ifpackageloaded{float}{}{\usepackage{float}}
\floatstyle{ruled}
\@ifundefined{c@chapter}{\newfloat{codelisting}{h}{lop}}{\newfloat{codelisting}{h}{lop}[chapter]}
\floatname{codelisting}{Listado}
\newcommand*\listoflistings{\listof{codelisting}{Listado de Listados}}
\makeatother
\makeatletter
\makeatother
\makeatletter
\@ifpackageloaded{caption}{}{\usepackage{caption}}
\@ifpackageloaded{subcaption}{}{\usepackage{subcaption}}
\makeatother
\makeatletter
\@ifpackageloaded{tikz}{}{\usepackage{tikz}}
\makeatother
        \newcommand*\circled[1]{\tikz[baseline=(char.base)]{
          \node[shape=circle,draw,inner sep=1pt] (char) {{\scriptsize#1}};}}  
                  

\usepackage{bookmark}
\IfFileExists{xurl.sty}{\usepackage{xurl}}{} % add URL line breaks if available
\urlstyle{same}
\hypersetup{
  pdftitle={Series de Pandas},
  pdfauthor={Mario Morales},
  pdflang={es},
  hidelinks,
  pdfcreator={LaTeX via pandoc}}


\title{Series de Pandas}
\subtitle{\emph{Estadística con Python}}
\author{Mario Morales}
\date{}

\begin{document}
\frame{\titlepage}

\renewcommand*\contentsname{Tabla de contenidos}
\begin{frame}[allowframebreaks]
  \frametitle{Tabla de contenidos}
  \setcounter{tocdepth}{3}
  \tableofcontents
\end{frame}
\setcounter{tocdepth}{3}
\tableofcontents
}

\begin{frame}[fragile]{Series de Pandas}
\phantomsection\label{series-de-pandas}
\begin{block}{Pandas}
\phantomsection\label{pandas}
\index{Pandas}

\begin{itemize}
\item
  \emph{Pandas} es una biblioteca de Python ampliamente utilizada para
  el análisis y procesamiento de datos.
\item
  Proporciona estructuras de datos rápidas, flexibles y útiles para
  cálculo estadístico y ofrece funciones que facilitan la entrada,
  organización y transformación de datos.
\item
  Las dos principales estructuras de datos de Pandas, \texttt{Series} y
  \texttt{DataFrame}, agotan la gran mayoría de los casos de uso típicos
  en estadística.
\end{itemize}
\end{block}

\begin{block}{Series de pandas}
\phantomsection\label{series-de-pandas-1}
\index{Diccionarios}

\begin{tcolorbox}[enhanced jigsaw, left=2mm, rightrule=.15mm, opacityback=0, title=\textcolor{quarto-callout-note-color}{\faInfo}\hspace{0.5em}{Nota}, toprule=.15mm, arc=.35mm, bottomrule=.15mm, colback=white, opacitybacktitle=0.6, leftrule=.75mm, bottomtitle=1mm, colbacktitle=quarto-callout-note-color!10!white, toptitle=1mm, colframe=quarto-callout-note-color-frame, titlerule=0mm, coltitle=black, breakable]

Una \emph{Serie} de Pandas es un arreglo unidimensional etiquetado
(\texttt{ndarray}). Este objeto, que admite la {indexación basada en
enteros y etiquetas}, proporciona una gran cantidad de métodos para
realizar operaciones relacionadas con el índice.

\end{tcolorbox}
\end{block}
\end{frame}

\begin{frame}[fragile]{Creación de Series}
\phantomsection\label{creaciuxf3n-de-series}
\begin{block}{Creación Series de Pandas}
\phantomsection\label{creaciuxf3n-series-de-pandas}
\index{Contenedores de datos!Series de Pandas} \index{Pandas!Series de}
\index{Series}

\begin{itemize}
\item
  A continuación se muestra como crear Series de Pandas.
\item
  Lo primero es importar el paquete \texttt{pandas}, lo hacemos con el
  nombre \texttt{pd} como es costumbre entre la comunidad de usuarios de
  esta biblioteca.
\end{itemize}

\begin{Shaded}
\begin{Highlighting}[]
\ImportTok{import}\NormalTok{ pandas }\ImportTok{as}\NormalTok{ pd}
\ImportTok{import}\NormalTok{ numpy }\ImportTok{as}\NormalTok{ np }
\NormalTok{datoS}\OperatorTok{=}\NormalTok{pd.Series([}\DecValTok{22}\NormalTok{,}\DecValTok{15}\NormalTok{,}\DecValTok{19}\NormalTok{]) }
\NormalTok{datoS}
\end{Highlighting}
\end{Shaded}

\begin{verbatim}
0    22
1    15
2    19
dtype: int64
\end{verbatim}

De manera predeterminada, la función \texttt{Series()} asigna índices
numéricos cuando, al momento de crear la serie, no se entregan las
etiquetas, como ocurre con los indices de un arreglo de \texttt{numpy} o
de una lista.
\end{block}

\begin{block}{Creación Series de Pandas (II)}
\phantomsection\label{creaciuxf3n-series-de-pandas-ii}
\begin{itemize}
\item
  Note que la salida impresa de la serie, además de los valores de las
  edades \texttt{22,\ 15\ \ y\ 19}, muestra las etiquetas (índices),
  asociadas a cada valor.
\item
  También podemos asignar las etiquetas, con el argumento
  \texttt{index}, al momento de crear la Serie, como se muestra a
  continuación:
\end{itemize}

\begin{Shaded}
\begin{Highlighting}[]
\NormalTok{datoS}\OperatorTok{=}\NormalTok{pd.Series([}\DecValTok{22}\NormalTok{,}\DecValTok{15}\NormalTok{,}\DecValTok{19}\NormalTok{],index}\OperatorTok{=}\NormalTok{[}\StringTok{\textquotesingle{}a\textquotesingle{}}\NormalTok{, }\StringTok{\textquotesingle{}b\textquotesingle{}}\NormalTok{, }\StringTok{\textquotesingle{}c\textquotesingle{}}\NormalTok{])}
\NormalTok{datoS}
\end{Highlighting}
\end{Shaded}

\begin{verbatim}
a    22
b    15
c    19
dtype: int64
\end{verbatim}
\end{block}

\begin{block}{Creación Series a partir de diccionarios}
\phantomsection\label{creaciuxf3n-series-a-partir-de-diccionarios}
\begin{itemize}
\item
  Es posible crear una serie de Pandas a partir de un diccionario.
\item
  Para ilustrarlo usaremos el diccionario \texttt{edad} creado en la
  clase.
\item
  Observe que las claves del diccionario pasan a ser las etiquetas de la
  serie.
\end{itemize}

\phantomsection\label{annotated-cell-3}%
\begin{Shaded}
\begin{Highlighting}[]
\NormalTok{edad}\OperatorTok{=}\NormalTok{\{}\StringTok{\textquotesingle{}Luis\textquotesingle{}}\NormalTok{:}\DecValTok{15}\NormalTok{,}\StringTok{\textquotesingle{}Brenda\textquotesingle{}}\NormalTok{:}\DecValTok{22}\NormalTok{,}\StringTok{\textquotesingle{}Diego\textquotesingle{}}\NormalTok{:}\DecValTok{19}\NormalTok{\}}
\NormalTok{edadS}\OperatorTok{=}\NormalTok{pd.Series(edad) }\hspace*{\fill}\NormalTok{\circled{1}}
\end{Highlighting}
\end{Shaded}

\begin{description}
\tightlist
\item[\circled{1}]
Serie creada a partir del diccionario.
\end{description}
\end{block}
\end{frame}

\begin{frame}[fragile]{Métodos para exploración de series}
\phantomsection\label{muxe9todos-para-exploraciuxf3n-de-series}
\begin{block}{Exploración inicial: métodos \texttt{.head} y
\texttt{.tail}}
\phantomsection\label{exploraciuxf3n-inicial-muxe9todos-.head-y-.tail}
\begin{itemize}
\item
  Cuando trabajamos con Series que contienen una gran cantidad de datos,
  suele ser poco práctico imprimirlos todos en pantalla.
\item
  Para una inspección rápida, usamos los métodos \texttt{.head(n)} y
  \texttt{.tail(n)}, los cuales permiten extraer los primeros y últimos
  \texttt{n} elementos de la Serie, respectivamente.
\item
  Si no se especifica el valor de \texttt{n}, por defecto ambos métodos
  devuelven 5 elementos.
\end{itemize}
\end{block}

\begin{block}{Ejemplo de uso de \texttt{.head} y \texttt{.tail}}
\phantomsection\label{ejemplo-de-uso-de-.head-y-.tail}
\begin{Shaded}
\begin{Highlighting}[]
\CommentTok{\# Creamos una serie más larga para el ejemplo}
\NormalTok{datos }\OperatorTok{=}\NormalTok{ pd.Series([}\DecValTok{12}\NormalTok{, }\DecValTok{45}\NormalTok{, }\DecValTok{32}\NormalTok{, }\DecValTok{18}\NormalTok{, }\DecValTok{56}\NormalTok{, }\DecValTok{29}\NormalTok{, }\DecValTok{78}\NormalTok{, }\DecValTok{91}\NormalTok{, }\DecValTok{14}\NormalTok{, }\DecValTok{63}\NormalTok{, }\DecValTok{88}\NormalTok{])}

\NormalTok{datos.head() }
\end{Highlighting}
\end{Shaded}

\begin{verbatim}
0    12
1    45
2    32
3    18
4    56
dtype: int64
\end{verbatim}

\begin{Shaded}
\begin{Highlighting}[]
\NormalTok{datos.tail(}\DecValTok{3}\NormalTok{) }
\end{Highlighting}
\end{Shaded}

\begin{verbatim}
8     14
9     63
10    88
dtype: int64
\end{verbatim}
\end{block}

\begin{block}{Identificación de valores únicos}
\phantomsection\label{identificaciuxf3n-de-valores-uxfanicos}
\index{Series!Método!unique}

\begin{itemize}
\item
  En muchas ocasiones, una Serie contiene datos categóricos o valores
  que se repiten múltiples veces, y nos interesa conocer únicamente
  cuáles son las distintas categorías que la componen.
\item
  El método \texttt{.unique()} recorre la Serie y retorna un arreglo de
  \texttt{NumPy} con los valores únicos, conservando el mismo orden en
  el que fueron encontrados por primera vez.
\end{itemize}
\end{block}

\begin{block}{Ejemplo de uso del método \texttt{.unique()}}
\phantomsection\label{ejemplo-de-uso-del-muxe9todo-.unique}
\begin{Shaded}
\begin{Highlighting}[]
\CommentTok{\# Ejemplo con categorías repetidas en un muestreo de laboratorio}
\NormalTok{especies }\OperatorTok{=}\NormalTok{ pd.Series([}\StringTok{\textquotesingle{}Cachama\textquotesingle{}}\NormalTok{, }\StringTok{\textquotesingle{}Tilapia\textquotesingle{}}\NormalTok{,}
              \StringTok{\textquotesingle{}Cachama\textquotesingle{}}\NormalTok{, }\StringTok{\textquotesingle{}Bocachico\textquotesingle{}}\NormalTok{,}
              \StringTok{\textquotesingle{}Cachama\textquotesingle{}}\NormalTok{, }\StringTok{\textquotesingle{}Tilapia\textquotesingle{}}\NormalTok{])}
\NormalTok{especies.unique() }
\end{Highlighting}
\end{Shaded}

\begin{verbatim}
array(['Cachama', 'Tilapia', 'Bocachico'], dtype=object)
\end{verbatim}
\end{block}
\end{frame}

\begin{frame}[fragile]{Recuperación de datos de una serie}
\phantomsection\label{recuperaciuxf3n-de-datos-de-una-serie}
\begin{block}{Recuperación de valores de la Serie}
\phantomsection\label{recuperaciuxf3n-de-valores-de-la-serie}
Como ocurre con los diccionarios, podemos recuperar los indices de una
serie mediante el método \texttt{.index} \index{Series!Método!index} y
los valores con el método \texttt{.values} \index{Series!Método!values}

\phantomsection\label{annotated-cell-7}%
\begin{Shaded}
\begin{Highlighting}[]
\NormalTok{edadS.index }\hspace*{\fill}\NormalTok{\circled{1}}
\NormalTok{edadS.values }\hspace*{\fill}\NormalTok{\circled{2}}
\end{Highlighting}
\end{Shaded}

\begin{description}
\tightlist
\item[\circled{1}]
Índices de la Serie.\\
\item[\circled{2}]
Valores de la Serie.
\end{description}

\begin{verbatim}
Index(['Luis', 'Brenda', 'Diego'], dtype='object')
array([15, 22, 19])
\end{verbatim}
\end{block}

\begin{block}{Recuperación de valores de la Serie}
\phantomsection\label{recuperaciuxf3n-de-valores-de-la-serie-1}
\begin{itemize}
\item
  Observe que los valores de la serie se recuperan como un arreglo de
  \texttt{NumPy}.
\item
  La recuperación de datos de una serie se hace de manera análoga a como
  se hace con los arreglos o los diccionarios

  \begin{itemize}
  \tightlist
  \item
    Por indice,
  \item
    Por etiqueta
  \item
    Por alguna condición que genere una lista, un arreglo o una Serie de
    tipo booleano
  \end{itemize}
\item
  Veamos algunos ejemplos.
\end{itemize}
\end{block}

\begin{block}{Recuperación de valores de la Serie: ejemplos}
\phantomsection\label{recuperaciuxf3n-de-valores-de-la-serie-ejemplos}
\phantomsection\label{annotated-cell-8}%
\begin{Shaded}
\begin{Highlighting}[]
\NormalTok{edadS[}\DecValTok{1}\NormalTok{] }\hspace*{\fill}\NormalTok{\circled{1}}
\NormalTok{edadS[}\StringTok{"Brenda"}\NormalTok{] }\hspace*{\fill}\NormalTok{\circled{2}}
\NormalTok{edadS[edadS}\OperatorTok{\textgreater{}}\DecValTok{20}\NormalTok{] }\hspace*{\fill}\NormalTok{\circled{3}}
\end{Highlighting}
\end{Shaded}

\begin{description}
\tightlist
\item[\circled{1}]
Segundo elemento de la Serie, seleccionado por índice. Si se entrega un
valor que esté por fuera del rango del indice, se produce un
\texttt{IndexError}.
\item[\circled{2}]
Equivalente al anterior, pero seleccionando por etiqueta. Si se entrega
una etiqueta que no esté en el índice se produce un \texttt{KeyError}
\item[\circled{3}]
Por la condición \texttt{edad} mayor que 20.
\texttt{edadS\textgreater{}20} produce una Serie con datos de tipo
booleano, \texttt{dtype(\textquotesingle{}bool\textquotesingle{})}, que
contiene \texttt{True} donde se cumpla la condición de ser mayor que 20
(\texttt{\textgreater{}20}) en \texttt{edadS}, y \texttt{False} en otro
caso.
\end{description}

\begin{verbatim}
22
22
Brenda    22
dtype: int64
\end{verbatim}
\end{block}

\begin{block}{Indexación: explícita vs.~implícita (\texttt{loc} e
\texttt{iloc})}
\phantomsection\label{indexaciuxf3n-expluxedcita-vs.-impluxedcita-loc-e-iloc}
\index{Series!Método!loc} \index{Series!Método!iloc}

\begin{itemize}
\tightlist
\item
  Aunque podemos extraer datos de una Serie usando corchetes
  \texttt{{[}{]}}, esto puede generar ambigüedad cuando nuestras
  etiquetas de índice también son números enteros.
\item
  Para evitar errores lógicos, Pandas provee dos atributos
  especializados:

  \begin{itemize}
  \tightlist
  \item
    \texttt{.loc{[}{]}}: Busca estrictamente por la \textbf{etiqueta
    explícita} del índice.
  \item
    \texttt{.iloc{[}{]}}: Busca por la \textbf{posición entera
    implícita} (desde 0 hasta n-1), de la misma forma que en las listas
    tradicionales.
  \end{itemize}
\end{itemize}
\end{block}

\begin{block}{Ejemplos de uso de \texttt{.loc} e \texttt{.iloc}}
\phantomsection\label{ejemplos-de-uso-de-.loc-e-.iloc}
\begin{Shaded}
\begin{Highlighting}[]
\NormalTok{edadS.loc[}\StringTok{\textquotesingle{}Diego\textquotesingle{}}\NormalTok{] }\CommentTok{\# loc si usas etiquetas }
\end{Highlighting}
\end{Shaded}

\begin{verbatim}
19
\end{verbatim}

\begin{Shaded}
\begin{Highlighting}[]
\NormalTok{edadS.iloc[}\DecValTok{2}\NormalTok{]   }\CommentTok{\# iloc si usas posición entera}
\end{Highlighting}
\end{Shaded}

\begin{verbatim}
19
\end{verbatim}

\begin{Shaded}
\begin{Highlighting}[]
\NormalTok{edadS.loc[}\DecValTok{2}\NormalTok{]  }
\NormalTok{edadS.iloc[}\StringTok{\textquotesingle{}Diego\textquotesingle{}}\NormalTok{]}
\end{Highlighting}
\end{Shaded}
\end{block}

\begin{block}{Ejemplos de uso de \texttt{.loc} e \texttt{.iloc}}
\phantomsection\label{ejemplos-de-uso-de-.loc-e-.iloc-1}
El escenario donde las etiquetas son números (como códigos de
estudiantes o identificadores, por ejemplo \texttt{101}, \texttt{102}),
hace obligatorio usar \texttt{.loc} e \texttt{.iloc} para evitar que el
programa se confunda entre la etiqueta y la posición.

\begin{Shaded}
\begin{Highlighting}[]
\CommentTok{\# Registro de biometría (peso en gramos) de especímenes de cachama en el CINPIC}
\NormalTok{peso}\OperatorTok{=}\NormalTok{ pd.Series([}\FloatTok{450.2}\NormalTok{, }\FloatTok{465.5}\NormalTok{, }\FloatTok{440.8}\NormalTok{], index}\OperatorTok{=}\NormalTok{[}\DecValTok{101}\NormalTok{, }\DecValTok{105}\NormalTok{, }\DecValTok{110}\NormalTok{])}

\NormalTok{peso.loc[}\DecValTok{105}\NormalTok{]  }\CommentTok{\# por etiqueta }
\end{Highlighting}
\end{Shaded}

\begin{verbatim}
465.5
\end{verbatim}

\begin{Shaded}
\begin{Highlighting}[]
\NormalTok{peso.iloc[}\DecValTok{1}\NormalTok{]   }\CommentTok{\# por posición }
\end{Highlighting}
\end{Shaded}

\begin{verbatim}
465.5
\end{verbatim}
\end{block}
\end{frame}

\begin{frame}[fragile]{Modificación del índice}
\phantomsection\label{modificaciuxf3n-del-uxedndice}
\begin{block}{Asignación futura del índice}
\phantomsection\label{asignaciuxf3n-futura-del-uxedndice}
\index{Series!Atributo!index}

\begin{itemize}
\item
  Es posible modificar o establecer las etiquetas de una Serie después
  de haberla creado, reasignando una nueva lista al atributo
  \texttt{.index}.
\item
  \textbf{¿Por qué es útil?} Es muy común que los datos numéricos
  provengan de un archivo o medición y carezcan de etiquetas iniciales .
  Posteriormente, obtenemos las etiquetas desde otra fuente y
  necesitamos integrarlas.
\item
  \textbf{Regla fundamental:} La cantidad de etiquetas nuevas
  proporcionadas debe ser exactamente igual a la cantidad de elementos
  en la Serie.
\end{itemize}
\end{block}

\begin{block}{Ejemplo de asignación del índice}
\phantomsection\label{ejemplo-de-asignaciuxf3n-del-uxedndice}
\phantomsection\label{annotated-cell-14}%
\begin{Shaded}
\begin{Highlighting}[]
\CommentTok{\# 1. Cargamos datos que inicialmente no tienen etiquetas descriptivas}
\NormalTok{temperaturas }\OperatorTok{=}\NormalTok{ pd.Series([}\FloatTok{28.5}\NormalTok{, }\FloatTok{30.2}\NormalTok{, }\FloatTok{29.8}\NormalTok{]) }
\NormalTok{temperaturas.index }

\CommentTok{\# 2. Obtenemos los códigos de las estaciones de monitoreo desde otra fuente}
\NormalTok{codigos\_estaciones }\OperatorTok{=}\NormalTok{ [}\StringTok{\textquotesingle{}Est\_Norte\textquotesingle{}}\NormalTok{, }\StringTok{\textquotesingle{}Est\_Centro\textquotesingle{}}\NormalTok{, }\StringTok{\textquotesingle{}Est\_Sur\textquotesingle{}}\NormalTok{]}

\CommentTok{\# 3. Asignamos los códigos como el nuevo índice de la Serie}
\NormalTok{temperaturas.index }\OperatorTok{=}\NormalTok{ codigos\_estaciones }\hspace*{\fill}\NormalTok{\circled{2}}

\NormalTok{temperaturas }\hspace*{\fill}\NormalTok{\circled{3}}
\end{Highlighting}
\end{Shaded}

\begin{verbatim}
RangeIndex(start=0, stop=3, step=1)
Est_Norte     28.5
Est_Centro    30.2
Est_Sur       29.8
dtype: float64
\end{verbatim}
\end{block}

\begin{block}{Índices con etiquetas duplicadas}
\phantomsection\label{uxedndices-con-etiquetas-duplicadas}
\index{Pandas!Índices duplicados}

\begin{itemize}
\tightlist
\item
  Una diferencia fundamental entre un diccionario estándar de Python y
  una Serie de Pandas, es que \textbf{Pandas sí permite tener etiquetas
  repetidas} en su índice.
\item
  Aunque para mantener los datos organizados se recomienda que los
  índices sean únicos, en la práctica es posible tener duplicados.
\item
  Si se utiliza \texttt{.loc{[}{]}} para extraer información mediante
  una etiqueta que está repetida, Pandas no devuelve un solo valor, sino
  una \textbf{nueva Serie} con todas las coincidencias.
\end{itemize}
\end{block}

\begin{block}{Ejemplo: etiquetas duplicadas}
\phantomsection\label{ejemplo-etiquetas-duplicadas}
\phantomsection\label{annotated-cell-15}%
\begin{Shaded}
\begin{Highlighting}[]
\CommentTok{\# Serie con la etiqueta \textquotesingle{}A\textquotesingle{} repetida}
\NormalTok{datos\_duplicados }\OperatorTok{=}\NormalTok{ pd.Series([}\DecValTok{10}\NormalTok{, }\DecValTok{20}\NormalTok{, }\DecValTok{30}\NormalTok{, }\DecValTok{40}\NormalTok{], index}\OperatorTok{=}\NormalTok{[}\StringTok{\textquotesingle{}A\textquotesingle{}}\NormalTok{, }\StringTok{\textquotesingle{}B\textquotesingle{}}\NormalTok{, }\StringTok{\textquotesingle{}A\textquotesingle{}}\NormalTok{, }\StringTok{\textquotesingle{}C\textquotesingle{}}\NormalTok{])}

\CommentTok{\# Búsqueda de una etiqueta única}
\NormalTok{datos\_duplicados.loc[}\StringTok{\textquotesingle{}B\textquotesingle{}}\NormalTok{] }\hspace*{\fill}\NormalTok{\circled{1}}

\CommentTok{\# Búsqueda de una etiqueta duplicada}
\NormalTok{datos\_duplicados.loc[}\StringTok{\textquotesingle{}A\textquotesingle{}}\NormalTok{] }\hspace*{\fill}\NormalTok{\circled{2}}
\end{Highlighting}
\end{Shaded}

\begin{verbatim}
20
A    10
A    30
dtype: int64
\end{verbatim}
\end{block}
\end{frame}

\begin{frame}[fragile]{Agregar datos a una serie}
\phantomsection\label{agregar-datos-a-una-serie}
\begin{block}{Ejemplo: agregar un elemento}
\phantomsection\label{ejemplo-agregar-un-elemento}
Agregamos un nuevo individuo a la serie \texttt{edadS}

\begin{Shaded}
\begin{Highlighting}[]
\NormalTok{edadS[}\StringTok{"Pedro"}\NormalTok{]}\OperatorTok{=}\DecValTok{16}

\NormalTok{edadS.tail(}\DecValTok{3}\NormalTok{)}
\end{Highlighting}
\end{Shaded}

\begin{verbatim}
Brenda    22
Diego     19
Pedro     16
dtype: int64
\end{verbatim}
\end{block}

\begin{block}{Agregar múltiples elementos a una Serie}
\phantomsection\label{agregar-muxfaltiples-elementos-a-una-serie}
Para agregar \textbf{múltiples} elementos nuevos al mismo tiempo,
debemos crear una nueva Serie con los datos a ingresar y luego unirlas
utilizando la función \texttt{pd.concat()}.

\begin{Shaded}
\begin{Highlighting}[]
\NormalTok{nuevos}\OperatorTok{=}\NormalTok{ pd.Series(\{}\StringTok{\textquotesingle{}Ana\textquotesingle{}}\NormalTok{: }\DecValTok{20}\NormalTok{, }\StringTok{\textquotesingle{}Juan\textquotesingle{}}\NormalTok{: }\DecValTok{25}\NormalTok{, }\StringTok{\textquotesingle{}Sofia\textquotesingle{}}\NormalTok{: }\DecValTok{21}\NormalTok{\})}

\CommentTok{\# Concatenamos ambas estructuras}
\NormalTok{edadS }\OperatorTok{=}\NormalTok{ pd.concat([edadS, nuevos])}
\NormalTok{edadS}
\end{Highlighting}
\end{Shaded}

\begin{verbatim}
Luis      15
Brenda    22
Diego     19
Pedro     16
Ana       20
Juan      25
Sofia     21
dtype: int64
\end{verbatim}
\end{block}

\begin{block}{Cuidado!!!}
\phantomsection\label{cuidado}
\begin{tcolorbox}[enhanced jigsaw, left=2mm, rightrule=.15mm, opacityback=0, title=\textcolor{quarto-callout-warning-color}{\faExclamationTriangle}\hspace{0.5em}{Advertencia}, toprule=.15mm, arc=.35mm, bottomrule=.15mm, colback=white, opacitybacktitle=0.6, leftrule=.75mm, bottomtitle=1mm, colbacktitle=quarto-callout-warning-color!10!white, toptitle=1mm, colframe=quarto-callout-warning-color-frame, titlerule=0mm, coltitle=black, breakable]

En versiones antiguas de Pandas existía el método \texttt{.append()},
que usamos en listas, pero fue eliminado a partir de la versión 2.0 por
razones de rendimiento.

\end{tcolorbox}
\end{block}
\end{frame}

\begin{frame}[fragile]{Conversión de tipos}
\phantomsection\label{conversiuxf3n-de-tipos}
\begin{block}{Modificación del tipo de dato}
\phantomsection\label{modificaciuxf3n-del-tipo-de-dato}
\index{Series!Método!astype}

\begin{itemize}
\tightlist
\item
  En ocasiones, los datos de una Serie no tienen el formato adecuado
  para el análisis. Por ejemplo, números que fueron importados por el
  sistema como texto (\texttt{string}).
\item
  El método \texttt{.astype()} permite forzar la conversión de los
  elementos de una Serie a un tipo de dato (\texttt{dtype}) específico,
  como enteros (\texttt{int}) o decimales (\texttt{float}).
\item
  Es importante recordar que este método devuelve una \textbf{nueva}
  Serie con los tipos cambiados; no modifica la original a menos que se
  sobreescriba.
\end{itemize}
\end{block}

\begin{block}{Ejemplo de modificación de tipo}
\phantomsection\label{ejemplo-de-modificaciuxf3n-de-tipo}
Serie donde los números están almacenados como texto

\begin{Shaded}
\begin{Highlighting}[]
\NormalTok{texto }\OperatorTok{=}\NormalTok{ pd.Series([}\StringTok{"15.5"}\NormalTok{, }\StringTok{"20.1"}\NormalTok{, }\StringTok{"18.3"}\NormalTok{])}
\NormalTok{texto.dtype }
\end{Highlighting}
\end{Shaded}

\begin{verbatim}
dtype('O')
\end{verbatim}

Conversión a números de punto flotante (decimales)

\phantomsection\label{annotated-cell-19}%
\begin{Shaded}
\begin{Highlighting}[]
\NormalTok{numerico }\OperatorTok{=}\NormalTok{ texto.astype(}\BuiltInTok{float}\NormalTok{) }\hspace*{\fill}\NormalTok{\circled{2}}
\NormalTok{numerico.dtype }\hspace*{\fill}\NormalTok{\circled{3}}
\end{Highlighting}
\end{Shaded}

\begin{verbatim}
dtype('float64')
\end{verbatim}
\end{block}
\end{frame}

\begin{frame}[fragile]{Alineación automática de índices}
\phantomsection\label{alineaciuxf3n-automuxe1tica-de-uxedndices}
\begin{block}{Alineación automática de índices en operaciones}
\phantomsection\label{alineaciuxf3n-automuxe1tica-de-uxedndices-en-operaciones}
\begin{itemize}
\tightlist
\item
  Una de las características más potentes de Pandas es que, al realizar
  operaciones matemáticas entre dos Series, los datos se \textbf{alinean
  automáticamente basándose en sus etiquetas}, no en su posición.
\item
  Si una etiqueta está presente en una Serie pero no en la otra, Pandas
  introduce un valor faltante (\texttt{NaN}) en el resultado para esa
  etiqueta.
\item
  Esto evita errores graves (que sí ocurrirían con listas normales)
  cuando se cruzan conjuntos de datos desordenados o con diferente
  cantidad de registros.
\end{itemize}
\end{block}

\begin{block}{Ejemplo: Alineación automática de índices}
\phantomsection\label{ejemplo-alineaciuxf3n-automuxe1tica-de-uxedndices}
\phantomsection\label{annotated-cell-20}%
\begin{Shaded}
\begin{Highlighting}[]
\CommentTok{\# Notas del primer y segundo corte de un grupo de estudiantes}
\NormalTok{corte1 }\OperatorTok{=}\NormalTok{ pd.Series(\{}\StringTok{\textquotesingle{}Luis\textquotesingle{}}\NormalTok{: }\FloatTok{4.0}\NormalTok{, }\StringTok{\textquotesingle{}Brenda\textquotesingle{}}\NormalTok{: }\FloatTok{4.5}\NormalTok{, }\StringTok{\textquotesingle{}Diego\textquotesingle{}}\NormalTok{: }\FloatTok{3.8}\NormalTok{, }\StringTok{\textquotesingle{}Ana\textquotesingle{}}\NormalTok{: }\FloatTok{4.2}\NormalTok{\})}
\NormalTok{corte2 }\OperatorTok{=}\NormalTok{ pd.Series(\{}\StringTok{\textquotesingle{}Brenda\textquotesingle{}}\NormalTok{: }\FloatTok{4.6}\NormalTok{, }\StringTok{\textquotesingle{}Ana\textquotesingle{}}\NormalTok{: }\FloatTok{4.0}\NormalTok{, }\StringTok{\textquotesingle{}Diego\textquotesingle{}}\NormalTok{: }\FloatTok{4.1}\NormalTok{, }\StringTok{\textquotesingle{}Pedro\textquotesingle{}}\NormalTok{: }\FloatTok{3.5}\NormalTok{\}) }

\CommentTok{\# Sumamos ambas series para obtener el acumulado}
\NormalTok{suma}\OperatorTok{=}\NormalTok{ corte1 }\OperatorTok{+}\NormalTok{ corte2 }\hspace*{\fill}\NormalTok{\circled{1}}

\NormalTok{suma }
\end{Highlighting}
\end{Shaded}

\begin{verbatim}
Ana       8.2
Brenda    9.1
Diego     7.9
Luis      NaN
Pedro     NaN
dtype: float64
\end{verbatim}
\end{block}

\begin{block}{Operaciones con índices duplicados}
\phantomsection\label{operaciones-con-uxedndices-duplicados}
\index{Pandas!Alineación con duplicados}

\begin{itemize}
\tightlist
\item
  El comportamiento matemático de Pandas con etiquetas repetidas depende
  de si las Series comparten exactamente la misma estructura de índices:

  \begin{itemize}
  \tightlist
  \item
    \textbf{Índices idénticos:} Pandas optimiza la operación y la
    realiza componente a componente.
  \item
    \textbf{Índices diferentes (con duplicados):} Pandas aplica un
    producto cartesiano cruzando todas las combinaciones posibles para
    las etiquetas en común.
  \end{itemize}
\end{itemize}
\end{block}

\begin{block}{Ejemplos con índice duplicado}
\phantomsection\label{ejemplos-con-uxedndice-duplicado}
\phantomsection\label{annotated-cell-21}%
\begin{Shaded}
\begin{Highlighting}[]
\CommentTok{\# Caso 1: Índices idénticos (Componente a componente)}
\NormalTok{serie1 }\OperatorTok{=}\NormalTok{ pd.Series([}\DecValTok{1}\NormalTok{, }\DecValTok{2}\NormalTok{], index}\OperatorTok{=}\NormalTok{[}\StringTok{\textquotesingle{}X\textquotesingle{}}\NormalTok{, }\StringTok{\textquotesingle{}X\textquotesingle{}}\NormalTok{])}
\NormalTok{serie2 }\OperatorTok{=}\NormalTok{ pd.Series([}\DecValTok{10}\NormalTok{, }\DecValTok{20}\NormalTok{], index}\OperatorTok{=}\NormalTok{[}\StringTok{\textquotesingle{}X\textquotesingle{}}\NormalTok{, }\StringTok{\textquotesingle{}X\textquotesingle{}}\NormalTok{])}
\NormalTok{serie1 }\OperatorTok{+}\NormalTok{ serie2 }\hspace*{\fill}\NormalTok{\circled{1}}

\CommentTok{\# Caso 2: Índices diferentes con duplicados (Producto cartesiano)}
\NormalTok{serie3 }\OperatorTok{=}\NormalTok{ pd.Series([}\DecValTok{1}\NormalTok{, }\DecValTok{2}\NormalTok{], index}\OperatorTok{=}\NormalTok{[}\StringTok{\textquotesingle{}X\textquotesingle{}}\NormalTok{, }\StringTok{\textquotesingle{}Y\textquotesingle{}}\NormalTok{])}
\NormalTok{serie4 }\OperatorTok{=}\NormalTok{ pd.Series([}\DecValTok{10}\NormalTok{, }\DecValTok{20}\NormalTok{], index}\OperatorTok{=}\NormalTok{[}\StringTok{\textquotesingle{}X\textquotesingle{}}\NormalTok{, }\StringTok{\textquotesingle{}X\textquotesingle{}}\NormalTok{])}
\NormalTok{serie3 }\OperatorTok{+}\NormalTok{ serie4 }\hspace*{\fill}\NormalTok{\circled{2}}
\end{Highlighting}
\end{Shaded}

\begin{verbatim}
X    11
X    22
dtype: int64
X    11.0
X    21.0
Y     NaN
dtype: float64
\end{verbatim}
\end{block}
\end{frame}

\begin{frame}[fragile]{Datos faltantes en una Serie}
\phantomsection\label{datos-faltantes-en-una-serie}
\begin{block}{Datos faltantes en series de Pandas}
\phantomsection\label{datos-faltantes-en-series-de-pandas}
\index{Pandas!Series!datos faltantes}

\index{Datos!faltantes}

\begin{itemize}
\item
  Las series de Pandas tienen la propiedad de poder tratar con datos
  faltantes,
\item
  Los cuales se representan como: \texttt{"NA"\ =\ Not\ available},
  \texttt{"NaN"\ =\ Not\ a\ number},
  \texttt{None\ =\ Null/nonetype\ object}.
\item
  A continuación se crean series que tienen datos faltantes
\end{itemize}
\end{block}

\begin{block}{Datos faltantes en series de Pandas}
\phantomsection\label{datos-faltantes-en-series-de-pandas-1}
\phantomsection\label{annotated-cell-22}%
\begin{Shaded}
\begin{Highlighting}[numbers=left,,]
\NormalTok{edadSN}\OperatorTok{=}\NormalTok{pd.Series([}\DecValTok{21}\NormalTok{, }\DecValTok{15}\NormalTok{, }\VariableTok{None}\NormalTok{], index}\OperatorTok{=}\NormalTok{[}\StringTok{\textquotesingle{}A\textquotesingle{}}\NormalTok{, }\StringTok{\textquotesingle{}B\textquotesingle{}}\NormalTok{, }\StringTok{\textquotesingle{}C\textquotesingle{}}\NormalTok{]) }\hspace*{\fill}\NormalTok{\circled{1}}
\NormalTok{edadSN}\OperatorTok{=}\NormalTok{pd.Series([}\DecValTok{21}\NormalTok{, }\DecValTok{15}\NormalTok{,np.nan],index}\OperatorTok{=}\NormalTok{[}\StringTok{\textquotesingle{}A\textquotesingle{}}\NormalTok{, }\StringTok{\textquotesingle{}B\textquotesingle{}}\NormalTok{, }\StringTok{\textquotesingle{}C\textquotesingle{}}\NormalTok{]) }\hspace*{\fill}\NormalTok{\circled{2}}
\NormalTok{edadDic}\OperatorTok{=}\NormalTok{\{}\StringTok{\textquotesingle{}A\textquotesingle{}}\NormalTok{:}\DecValTok{12}\NormalTok{, }\StringTok{\textquotesingle{}B\textquotesingle{}}\NormalTok{:}\DecValTok{15}\NormalTok{, }\StringTok{\textquotesingle{}C\textquotesingle{}}\NormalTok{:}\VariableTok{None}\NormalTok{\} }\hspace*{\fill}\NormalTok{\circled{3}}
\NormalTok{edadSN}\OperatorTok{=}\NormalTok{pd.Series(edadDic) }\hspace*{\fill}\NormalTok{\circled{4}}
\NormalTok{edadSN}\OperatorTok{=}\NormalTok{pd.Series(edadDic,index}\OperatorTok{=}\NormalTok{[}\StringTok{"A"}\NormalTok{,}\StringTok{"B"}\NormalTok{,}\StringTok{"D"}\NormalTok{]) }\hspace*{\fill}\NormalTok{\circled{5}}
\NormalTok{edadSN}
\end{Highlighting}
\end{Shaded}

\begin{description}
\tightlist
\item[\circled{1}]
Serie creada a partir de una lista con un valor en \texttt{None}.
\item[\circled{2}]
Serie creada a partir de un arreglo de Numpy con un valor en
\texttt{nan}.
\item[\circled{3}]
Diccionario con una clave cuyo valor es \texttt{None}
\item[\circled{4}]
Serie creada a partir del diccionario, el valor asociado al índice
\texttt{C} es \texttt{NaN}.
\item[\circled{5}]
\texttt{D} no es una clave del diccionario \texttt{edadDic}, aparece
como índice en la Serie con valor \texttt{NaN}.
\end{description}

\begin{verbatim}
A    12.0
B    15.0
D     NaN
dtype: float64
\end{verbatim}
\end{block}

\begin{block}{Datos faltantes en series de Pandas}
\phantomsection\label{datos-faltantes-en-series-de-pandas-2}
\begin{itemize}
\item
  En la penúltima linea del bloque anterior, observe que la clave
  \texttt{C} es descartada y no aparece en la Serie creada, a pesar de
  encontrarse en el diccionario, eso ocurre porque \texttt{C} no aparece
  en el argumento pasado como índice.
\item
  En contraste, \texttt{D} no está como clave en el diccionario, pero
  está en la lista que se pasó como índice al crear la Serie, por eso se
  le asigna el valor \texttt{NaN}.
\end{itemize}

\index{Series!Método!notnull()}
\end{block}

\begin{block}{¿Como saber si hay datos faltantes?}
\phantomsection\label{como-saber-si-hay-datos-faltantes}
\begin{itemize}
\item
  Los métodos \texttt{isnull()} y \texttt{notnull()} nos permiten
  conocer si en una serie de Pandas hay datos faltantes.
\item
  El método \texttt{isnull()}, aplicado a una Serie, regresa una Serie
  booleana con valores \texttt{True} donde hayan datos faltantes, y
  \texttt{False} en caso contrario.
\item
  El método \texttt{notnull()} regresa la negación lógica de
  \texttt{isnull()}, es decir, regresa valores \texttt{True} donde no
  hayan datos faltantes, y \texttt{False} en caso contrario.
  \index{Series!Método!all}
\item
  El método \texttt{all()} aplicable a una serie de Pandas que contiene
  valores booleanos (\texttt{dtype:\ bool}) regresa \texttt{True} en
  caso que todos los elementos de la serie sean \texttt{True} y
  \texttt{False} en caso que al menos uno elemento sea \texttt{False}.
\item
  La combinación de \texttt{isnull()} o \texttt{notnull()} con
  \texttt{all()} es útil para determinar si en una serie hay o no datos
  faltantes. En las siguientes líneas de código, se ilustra el
  funcionamiento de estos métodos.
\end{itemize}
\end{block}

\begin{block}{}
\phantomsection\label{section}
\phantomsection\label{annotated-cell-23}%
\begin{Shaded}
\begin{Highlighting}[]
\NormalTok{edadSN.isnull() }\hspace*{\fill}\NormalTok{\circled{1}}
\NormalTok{edadSN.notnull() }\hspace*{\fill}\NormalTok{\circled{2}}
\NormalTok{edadSN.notnull().}\BuiltInTok{all}\NormalTok{() }\hspace*{\fill}\NormalTok{\circled{3}}
\end{Highlighting}
\end{Shaded}

\begin{description}
\tightlist
\item[\circled{1}]
Serie de valores lógicos, con \texttt{True} en la posición donde hay un
dato faltante.
\item[\circled{2}]
Serie de valores lógicos, con \texttt{True} en caso que no hayan datos
faltantes, \texttt{False} en caso que haya al menos un dato faltante.
\item[\circled{3}]
El resultado \texttt{False} nos permite concluir que en la serie hay por
lo menos un dato faltante.
\end{description}

\begin{verbatim}
A    False
B    False
D     True
dtype: bool
A     True
B     True
D    False
dtype: bool
False
\end{verbatim}
\end{block}

\begin{block}{Tratamiento de datos faltantes en Series}
\phantomsection\label{tratamiento-de-datos-faltantes-en-series}
\begin{itemize}
\item
  Hay dos formas de proceder cuando se tienen datos faltantes:

  \begin{itemize}
  \tightlist
  \item
    Una es simplemente eliminarlos, lo cual se hace, en Series de
    Pandas, con el método \texttt{.dropna()}.
  \item
    La otra se conoce como \index{Imputación}
    \index{Datos faltantes!Imputación} \emph{imputación}, que consiste
    en estimar el (los) dato(s) faltante(s) mediante algún
    procedimiento.
  \end{itemize}
\item
  Se cuenta con mucha literatura en cuanto a métodos de imputación, uno
  sencillo es el de medias, que consiste en reemplazar el dato faltante
  por el promedio de los datos disponibles, también se puede usar la
  mediana.
\end{itemize}
\end{block}

\begin{block}{Imputación de datos faltantes en series de Pandas}
\phantomsection\label{imputaciuxf3n-de-datos-faltantes-en-series-de-pandas}
\begin{itemize}
\item
  A continuación se \emph{imputa} el dato faltante en la serie
  \texttt{edadSN} reemplazándolo por la media de los datos disponibles,
  la cual se calcula con \texttt{edadSN.mean()}.
\item
  De manera análoga \index{Pandas!Datos faltantes!Imputación} se hace
  imputación de datos faltantes usando la mediana, calculada con
  \texttt{edadSN.median()}.
\item
  En clases posteriores studiaremos con más detalles estos métodos,
  aplicables a Series de Pandas, para cálculos de estadísticas
  descriptivas.
\end{itemize}
\end{block}

\begin{block}{Imputación de datos faltantes (cont)}
\phantomsection\label{imputaciuxf3n-de-datos-faltantes-cont}
En las siguientes líneas de código, primero eliminamos los registros
faltantes y luego imputamos

\phantomsection\label{annotated-cell-24}%
\begin{Shaded}
\begin{Highlighting}[]
\NormalTok{edadSN.dropna() }\hspace*{\fill}\NormalTok{\circled{1}}
\NormalTok{edadSN.fillna(edadSN.mean()) }\hspace*{\fill}\NormalTok{\circled{2}}
\NormalTok{edadSN.fillna(edadSN.median()) }\hspace*{\fill}\NormalTok{\circled{3}}
\end{Highlighting}
\end{Shaded}

\begin{description}
\tightlist
\item[\circled{1}]
Elimina de la Serie los índices con datos faltantes.
\item[\circled{2}]
Reemplaza en la Serie los datos faltantes por la \textbf{media} de los
datos disponibles.\\
\item[\circled{3}]
Reemplaza en la Serie los datos faltantes por la \textbf{mediana} de los
datos disponibles.
\end{description}

\begin{verbatim}
A    12.0
B    15.0
dtype: float64
A    12.0
B    15.0
D    13.5
dtype: float64
A    12.0
B    15.0
D    13.5
dtype: float64
\end{verbatim}
\end{block}

\begin{block}{}
\phantomsection\label{section-1}
Tenga presente que en ninguna de las tres líneas de código anteriores se
modifica la Serie \texttt{edadSN}, si quiere conservar los resultados
deberá asignarlos a otro objeto o reemplazar, por ejemplo,
\texttt{edadSN=edadSN.dropna()}.
\end{block}
\end{frame}




\end{document}
